\section{Problem Setting}

\subsection{Problem Formulation}
%\biplav{Use from \cite{srivastava2022ultra} paper.}

In this section, we introduce the domain of our work and provide the groundwork for the problem. Funding agencies periodically announce Requests For Proposals (RFPs) on specific themes where they are looking for ideas to fund. Researchers in turn respond to those calls, with proposals, where they explain their ideas and detail a meaningful and actionable plan for how they plan to achieve said goals within a specific allotted budget, time frame, and other constraints. In doing so, they also often look to team with other colleagues to respond to such calls. As such, we consider a teaming environment where the availability of candidate participants may change at any given time, often sporadically.

In terms of terminology, let $S$ denote the set of all skills. Then, the {\em demand} for teaming is represented by the set of skills $S_i$ desired by an RFP $c_i$. Further, the  {\em supply}  is represented at an institution by the set of researchers $R$ along with their respective research interests and profiles to satisfy the demand.  
Teaming objectives may be short-term ($ST$), long-term ($LT$), or a combination of the two. An example of a short-term goal would be to meet the immediate requirements and skill needs, $S=\{s_1, s_2, ..., s_\alpha\}$ set forth by a call for proposal $c_i$. Each call in $C=\{c_1, c_2, ..., c_M\}$ requires a very specific skill set $S_i$ and an immediate and optimal assembling of available candidate researchers, $R=\{r_1, r_2, ..., r_N\}$. We then solve our teaming setup in three phases:~(1)~we first \textit{match} all researchers who may be of interest to the calls based on skills needed,~(2)~\textit{group} which subset of researchers should be recommended to be in a team $t_i$, and~(3)~compute goodness score $g_i$ for each team in $T=\{t_1, t_2, ..., t_\beta\}$ and recommend the top-\textit{k} suggestions to interested users. 

Similarly, some example scenarios of a long-term objective would be to maximize the number of funded awards $A$ %$A=\{a_1, a_2, ..., a_\gamma\}$ 
given to a researcher $r_j$ over a  time period (${LT}^A_t$), have a robust (diversified) pool of experienced talent ($LT^R_t$), and satisfy diversity goals of researchers' institutions.

%The team building setup that we consider is shown in Figure~\ref{fig:teaming_setup}. We have two types of interested users for our system: 
Our system would be of interest to at least two types of users: (1)~administrators at researchers’ organizations (e.g., university institutions) who want to promote more collaborations, proposals, and diversity at their institutions, and (2)~potential team members who will respond to a given RFP $c_i$ and are looking for collaborative opportunities. %The top left of the figure shows that 
Various environments call for different teams to be formed and matched with relevant opportunities. The candidate member set also may change over time, along with each of their skills and research interests. 

%\biplav{Talk about skills (e.g., Ultra-metric - s1, s2, s3, etc. Introduce terminology for the rest of the paper here. "Ci" for call for proposals, "Ri" for set of researchers, "Ti" for teams, "M" for number of calls, "N" number of researchers, "somethingsomething" for number of teams, "alpha" for set of skills 1-), "fi(C\_i)" gives skills for a RFP, "fi(R\_i)" gives skills for a researcher, "sigma(string)" gives ACM classification code (method-specific), team building formulation/equation}

Each user (i.e., admin or researcher) will interact with the system when a call for proposal $c_i$ is announced. Based on the requirements set by $c_i$, along with profiles extracted from researchers, the system will then algorithmically suggest teaming choices $T=\{t_1, t_2, ..., t_\beta\}$, which the users may accept or reject. We evaluate the efficacy of our algorithms and validate the teaming outputs via a goodness score. %(see Section~\ref{subsec:metrics}).

%\vspace{-0.15in}
\subsection{Use Cases} \label{subsec:usecases}
We show at least three practical use cases (UCs) to demonstrate the utility of our system. Each use case has various input prompts to select from and includes different algorithms that can be used to recommend or suggest proposals and teams to interested users. 

% \begin{figure}
%     \centering
%     \includegraphics[width=\linewidth]{cikm2023-applied/figures/ultra_system_uc1.png}
%     \caption{${UC}_1$ for ULTRA: Given a researcher $r_j$ and matching method, display calls for proposals and teaming suggestions.}
%     \label{fig:ultra_system_usecases}
% \end{figure}

For the first use case, ${UC}_1$, given a researcher's name and a teaming method $M_i$, we display a list of $k$ highest ranked proposals and possible teams (shown in Figure~\ref{fig:ultra_system_usecases}). Similarly, for ${UC}_2$, given a proposal call $c_i$ from a list of recently announced proposals $C$ (ideally refreshed in real-time or regularly), we display the best possible teams $T$ available for $c_i$. And the final use case, ${UC}_3$, takes input in the form of a research interest and teaming method, and displays respective matching proposals and teams for those parameters. 
%We show  screenshots of all the use cases in GitHub \cite{ultra-resources}. 
%Due to space constraints, we show only a truncated version of our UI (${UC}_1$) in Figure~\ref{fig:ultra_system_usecases}.
%Due to space constraints, we show only a truncated version of our UI (${UC}_1$) in Figure~\ref{fig:ultra_system_usecases}. 
We empirically evaluate the three use cases and the functionality of the four methods used in the later sections of the paper. %in the later sections of Section~\ref{subsec:humanstudy}. 